\section{Lean formalization: a concrete completion route}\label{sec:lean}

\subsection{What can be reused from ProveIt}

The repository was inspected at the immutable commit
\[
 \texttt{620f5e3b959c3a4284060ada5efb651812612319}.
\]
Its root toolchain specifies Lean \texttt{v4.32.0} and the
matching Mathlib tag. The relevant existing development is
\texttt{Logic/PrimitiveRecursiveArithmetic/ConciseIPRA}.
The following are actual declarations in the inspected sources
\cite{ProveIt}, not proposed names.

\begin{center}
\begin{tabular}{@{}p{0.38\textwidth}p{0.54\textwidth}@{}}
\toprule
Existing interface & Use in this project\\
\midrule
\code{PR}, \code{PTm} in \code{PR.lean} &
Primitive recursive function symbols and first-order terms.
The constructors of \code{PR} are zero, successor,
projection, composition, and primitive recursion.\\[3pt]
\code{PRAAx} &
Defining equations, successor nonzero and injectivity,
and the quantifier-free induction schema.\\[3pt]
\code{PRA.qfInd} &
Build an object-level induction derivation from a QF formula
and derivations of its base and successor cases.\\[3pt]
\code{Prf} in \code{Deduction.lean} &
The intuitionistic first-order derivation type.\\[3pt]
\code{NFm.ofPRA}, \code{pra\_iff\_native} &
Translate a completed iPRA derivation into the native
ConciseIPRA calculus.\\
\bottomrule
\end{tabular}
\end{center}

The repository's native/scalar equivalence does not automatically
turn a Mathlib theorem about real curves into an iPRA derivation.
Its use here is precise: after constructing the arithmetic
derivation, apply the existing theorem
\[
 (\Gamma\vdash[\mathsf{PRAAx}]\varphi)
   \ \longleftrightarrow\
 \mathsf{NPrf}\,(\Gamma.\mathsf{map}\,\mathsf{NFm.ofPRA})\,
                   (\mathsf{NFm.ofPRA}\,\varphi).
\]
The accompanying \code{NativeTransfer.lean} contains exactly
this generic wrapper with the real repository API.

The supplied draft directory was inspected for context. It contains
other constructive-analysis archives, but this proof needs none
of their claimed analytic results. In particular there is no
dependency on a formalized fundamental theorem of algebra,
an integral winding number, or an existing Jordan theorem.

\subsection{Three deliverable layers}

The completion should keep three layers distinct.

\begin{enumerate}
\item \textbf{Executable finite mathematics.}
Define rational points, half-open ray parity, meshes, distance
certificates, the finite honeycomb, reachability, boundary
traversal, and the numerical procedures from the bound ledger.
The Python implementation is a reference for this layer.
\item \textbf{Lean semantic verification.}
Prove the finite algorithms correct for their explicit data
structures. Named-curve theorems can then be stated using
fast Cauchy names, supplied moduli, and finite certificates.
This checks the mathematical construction.
\item \textbf{Object-level iPRA derivations.}
Represent the finite routines by \code{PR} expressions,
their Boolean invariants by equations of \code{PTm}, and
their proofs by \code{Prf PRAAx}. Use \code{PRA.qfInd}
for every loop invariant. Finally translate with
\code{pra\_iff\_native}.
\end{enumerate}

A proof in layer 2 alone is not a substitute for layer 3.
Conversely, once layer 3 produces a derivation in the explicit
intuitionistic calculus, a general analytic theorem from Mathlib
is unnecessary as a trusted premise.

\subsection{Module order and exact proof obligations}

\begin{longtable}{@{}p{0.20\textwidth}p{0.73\textwidth}@{}}
\toprule
Proposed module & Required results\\
\midrule
\endhead
\code{Coding} &
Integer/rational/list codes, bounded folds and scans, finite
minimum, powers and moduli composition. Supply \code{PR}
terms and defining-equation proofs.\\[3pt]
\code{Geometry} &
Rational segment membership and intersections; the
at-most-eight-candidate maximum-norm distance algorithm;
convex-combination norm inequalities.\\[3pt]
\code{Parity} &
The half-open edge test, xor additivity, reversal,
triangle parity outside its hull, convex-square and
parallelogram values. All cases are rational sign cases.\\[3pt]
\code{Replacement} &
Finite triangle-boundary cancellation, query-translation
parallelograms, mesh refinement, and stable curve parity
with explicit rational-margin certificates.\\[3pt]
\code{CycleWalk} &
Short-arc selection in a cyclic array; usage-bit incidence
at each vertex; constancy of all edge-usage bits;
Lemma~\ref{lem:factor} with the bound $sN$.\\[3pt]
\code{Honeycomb} &
The six vertices and neighbor table, vertical-slice
coverage, shared-edge identities, three-cell vertices,
and the center-link crossing and detour calculations.\\[3pt]
\code{FiniteGraph} &
$M$-round reachability, parent paths, stabilization,
monochromatic outer ring, degree-two cycle extraction
with at most $6M$ steps.\\[3pt]
\code{Connection} &
Near/safe link certificates; cycle proximity; local
parity jump; product-bit contradiction; route validity
and clearance at least $h$.\\[3pt]
\code{BothSides} &
The local square, half-arc nets, color certificates,
transition clearance, and the finite chain contradiction.
Produce rational witnesses with clearance $2^{-K-4}$.\\[3pt]
\code{Names} &
Cauchy names, padded moduli, finite distance nets,
completed evaluation, and the explicit circle-name lift.
Universal estimates are derived after finite certificates.\\[3pt]
\code{JordanSchema} &
Classification, route theorem, both-sides theorem,
size bounds, approximate intersection, and the
rational-witness common-boundary characterization.\\[3pt]
\code{NativeJordan} &
Assemble the closed arithmetic formulas and object-level
proofs; apply the existing native translation.\\
\bottomrule
\end{longtable}

The practical first milestone is \code{Parity} plus
\code{FiniteGraph}. It isolates the finite proof without
introducing a real-number library. The next milestone is
the factorization and routing theorem for finite sample
certificates. Only afterward should the names and moduli
interface connect those certificates to continuous curves.

There is no unknown topological lemma left in this plan.
The substantial work is formalizing the stated rational sign
lemmas, array invariants, and their object-theory derivations.

\subsection{Suggested statement and derivation interfaces}

The semantic layer should make certificates explicit. Suitable
structures include the following fields (this is an interface
specification, not a claim that these complete structures have
already been implemented).

\begin{description}[leftmargin=0pt,labelindent=0pt]
\item[\texttt{CurvePresentation}]
Fixed rational approximation routines, $U,V$, and the
arithmetic Cauchy/continuity/separation predicates.
\item[\texttt{OffCertificate}]
A dyadic clearance exponent and the corresponding
universal rational-name inequalities.
\item[\texttt{NetCertificate}]
The finite sample list, its parameter mesh and evaluation
precision, and the rational constants in its error bound.
\item[\texttt{RouteCertificate}]
Finite rational intermediate vertices, parent indices,
safe-link checks against a certified net, and a dyadic
clearance exponent.
\item[\texttt{BoundaryWitnesses}]
Two rational points, their clearance exponents, their
finite bit calculations, and finite precision witnesses
to the requested proximity.
\end{description}

The statement for routing should return actual data, for example
a route certificate satisfying the endpoint and clearance
predicates, rather than only a proposition asserting that two
points are in one connected component. The approximate crossing
statement should return the two rational sample parameters.
The reverse-boundary statement should return the displayed
primitive recursive circle-name construction.

For object-level proof assembly, the recurring pattern is:
\begin{enumerate}
\item Express a finite invariant as a Boolean equation
$I(s,\vec z)=1$, where $I$ is a fixed \code{PR} expression.
\item Prove its zero equation from the defining equations.
\item From $I(s,\vec z)=1$, prove $I(s+1,\vec z)=1$ using
the rational sign lemmas and Boolean equations.
\item Apply \code{PRA.qfInd}; only then introduce the
desired universal quantifiers or replace a Boolean fold
by its bounded mathematical reading.
\end{enumerate}
The final target is an actual object-level value of the schematic
type
\begin{lstlisting}
[] |- [PRAAx] jordanSentence
\end{lstlisting}
with the repository's notation and all the primitive recursive
witness terms substituted. Here \code{jordanSentence} is a
project-specific encoded sentence to be defined; the line is
not a completed declaration in this archive.
The supplied native wrapper takes such a derivation as an input.

The semantic definitions and the \code{PR} expressions should
be connected by evaluation lemmas. This prevents a discrepancy
between the executable reference algorithm and the routine
whose correctness is proved inside iPRA. A small library of
verified constructors for bounded scans and array iterations
will amortize most of the proof-term bookkeeping.

\subsection{Verification gates and honest status}

For a completed formalization:
\begin{enumerate}
\item Build against the pinned repository toolchain, then run
\code{lake env lean} on each module and the library build.
\item Check the semantic finite lemmas on the same exact
fixtures as the Python reference.
\item Audit every induction formula used in the object
derivations and inspect the generated \code{PR} expressions.
\item Inspect \code{\#print axioms} for the final semantic
and proof-construction theorems, and inspect their imports.
An object-level derivation of the specified arithmetic
sentence is the final iPRA gate.
\end{enumerate}

The environment used for this package did not contain Lean.
The files \code{lean/FiniteKernel.lean} and
\code{lean/NativeTransfer.lean} are therefore explicitly marked
\emph{uncompiled scaffolds}. They contain executable definitions
and a generic existing-API transfer wrapper, no added
\code{axiom} declarations and no \code{sorry} proofs.
They do not state that Jordan's theorem has already been
machine-checked. The complete dependency and invariant plan
is also in \code{lean/ROADMAP.md}.

\section{Executable artifacts and reproducibility}\label{sec:artifacts}

The archive includes an exact rational finite kernel, tests, and
machine-readable results. Its geometry uses Python's
\code{fractions.Fraction}; no floating-point tolerances occur in
the assertions. Plotting converts coordinates to floating point
only after the exact checks.

The kernel implements:
\begin{itemize}
\item the half-open ray index and its parity;
\item exact rational segment intersections and maximum-norm
clearances;
\item the rational honeycomb and bounded flood fill;
\item finite parent routes and separating cycles;
\item cyclic-mesh labeling and both integer and mod-two
factorization checks.
\end{itemize}
The signed integer index is an optional stronger helper in the
reference program. The proof in this article only trusts the
parity calculation.

The delivered run passed \textbf{142 exact checks}. Its concrete
connection examples are:

\begin{center}
\begin{tabular}{@{}p{0.47\textwidth}p{0.17\textwidth}rr@{}}
\toprule
Example & Outcome & Segments & Clearance\\
\midrule
Square: inside to inside & route & 3 & $11/8$\\
Square: inside to outside & bits $1,0$ & -- & --\\
Concave U: one arm to the other & route & 24 & $1/8$\\
Concave U: notch to far exterior & route & 25 & $3/16$\\
\bottomrule
\end{tabular}
\end{center}

The square separator has 106 edges. An additional annular
cell-coloring fixture produces two separating cycles, with
54 and 18 edges. This checks a particularly relevant case:
the finite boundary need not be a single cycle.
Other assertions cover vertex-height degeneracies, collinear
subdivision, reversal, repeated traversal, affine coordinate
changes, triangle replacement, all six neighbor identities,
and cyclic-mesh parity multiplication.

These checks validate the executable finite examples. They
are not an exhaustive check of the general theorem or a
substitute for its proof, and the reference implementation
uses polygons as finite distance oracles. The paper supplies
the certified net replacement for a general curve. The
Python program is not advertised as a completed extraction
of every continuous-name procedure in the paper.

To reproduce the finite checks and figures from the archive root:
\begin{lstlisting}
python3 artifacts/test_kernel.py
python3 artifacts/make_figures.py
\end{lstlisting}
The tests use only the Python standard library. The figure
script additionally needs Matplotlib. To compile the article:
\begin{lstlisting}
latexmk -pdf -interaction=nonstopmode -halt-on-error jordan_ipra.tex
\end{lstlisting}
All textual proof sections and the vector figure are included.
The README records the file layout and the verification
status in one place.
